\scanpage{64}
% archive sheet number: 63
sont de la forme $L(E;E)$, où $E$ est un fibré vectoriel de dim. $n$.

Si $\unsure{A|\mathscr{A}},\unsure{B|\mathscr{B}}$ sont deux fibrés en algèbres de \deleted{rang} $m^2,n^2$, alors $A\otimes B$ est un fibré en algèbres de rang $m^2n^2=(mn)^2$, donc définit un élément de $H^1(X,\underline{\operatorname{GP}}(mn-1))$. \unsure{Ceci} est associé à une représentation de groupes
\[
\operatorname{GP}(m-1)\times\operatorname{GP}(n-1)\unsure{-|\to}\operatorname{GP}(mn-1)
\]
déduite de
\[
\operatorname{Gl}(m)\times\operatorname{Gl}(n)\longrightarrow\operatorname{Gl}(mn)
\]
par passage au quotient. Ainsi l'ens
\[
\unsure{\mathscr{B}|\mathcal{B}|B}(X)=\coprod_{n\geq\unsure{1|2}}H^1(X,\operatorname{GP}(n-1))
\]
est un monoïde \inserted{comm., \deleted{associatif}}, dont l'élémnt unité est l'unique élément de degré $1$. On \deleted{\unsure{met}} introduit dans ce monoïde une relation d'équivalence : deux \deleted{\unsure{fibrés}} élémts $a,b$ sont équivalents, si on peut trouver $c$, $c$ étant ``trivial'' --- i.e. provenant d'un fibré vectoriel --- t.q. $ac=bc$. Si $a\sim b$ et $a'\sim b'$, alors $aa'\sim bb'$, donc l'ensemble quotient, noté $\operatorname{Br}(X)$, est un monoïde \deleted{C'est un} (ass, comm. avec $1$). C'est \unsure{même} un groupe, car si
